

\addcontentsline{toc}{chapter}{Abstract (Hebrew)}
\markboth{Abstract (Hebrew)}{Abstract (Hebrew)}

\chapter*{\texthebrew{תקציר}}
\label{sec:abstract-hebrew}
\begin{otherlanguage}{hebrew}
\begin{justify}
\begin{flushright}
% \textbf{ארכיטקטורות למידה עמוקה לתקשורת ממסר דו-קפיצתית: מחקר השוואתי של אסטרטגיות ממסר קלאסיות ומבוססות רשתות נוירונים}

% חיבור זה מציג מחקר השוואתי של אסטרטגיות ממסר (relay) קלאסיות ומבוססות רשתות נוירונים עבור תקשורת שיתופית דו-קפיצתית. המחקר כולו מתבצע על תצורה קנונית אחת, הקבועה מראש ואינה משתנה לאורך קו הטיעון המרכזי: קישור SISO בשני הקפיצות, מעל דעיכת Rayleigh מהירה בלתי-תלויה ושווה-התפלגות, בבסיס-רוחב מרוכב, אפנון BPSK, ושיעור שגיאות סיביות (BER) בלתי-מקודד, כאשר פונקציית הממסר היא הציר היחיד המשתנה. נבדקות תשע שיטות ממסר: שתי גישות קלאסיות : הגברה-והעברה (AF) ופענוח-והעברה (DF) : ושבע שיטות מבוססות למידה עמוקה: למידה מפוקחת (MLP, ממסר Hybrid המתאים את עצמו לרמת ה-SNR), מודלים גנרטיביים (VAE, CGAN), וארכיטקטורות רצפים (Transformer, Mamba S6, Mamba-2 SSD). כל התוצאות מבוססות על סימולציית מונטה קרלו (100,000 ביטים לכל נקודת SNR, עשרה ניסויים, רווחי סמך של 95\%).

% בתצורה הקנונית, ממסרי הרשת הטובים ביותר עולים על AF ב-SNR נמוך ($0$–$4$ dB) אך רק משתווים ל-DF הסימבולי, אשר שולט בתחום SNR בינוני-גבוה עם אפס פרמטרים. מחקר מורכבות חושף יחס U-הפוך בין גודל המודל לביצועים: רשת מינימלית בת 169 פרמטרים משתווה למודלים גדולים פי 100, בעוד גדלים בינוניים מציגים התאמת-יתר; בתקציב שווה של כ-3,000 פרמטרים הפער בין הארכיטקטורות כמעט נעלם, ממצא המעיד שהקיבולת, ולא הארכיטקטורה, היא הגורם המכריע במשימה זו.

% תרומה משנית מצומצמת בוחנת אפנונים מסדר גבוה: הממצאים הקנוניים עוברים במלואם ל-QPSK מכוח אי-תלות רכיבי I/Q, בעוד 16-QAM חושף רצפת BER מבנית של עיבוד מפוצל לפי ציר; ניסוח הממסר כמסווג דו-ממדי משותף על פני 16 נקודות הקונסטלציה מבטל רצפה זו ומשתווה ל-DF ב-SNR גבוה.

% התרומה המרכזית היא מחקר שיטתי של ממסר מבוסס-למידה תחת ערוצים \emph{בלתי-ידועים ובלתי-מותאמים}. בערוץ ISI בעל שלושה מקדמים, הבלתי-ידוע לשרשרת הקליטה, הממסרים הקלאסיים חסרי-הזיכרון ננעלים על רצפת BER אנליטית של 0.25 : כאשר ביצועי DF אף מחמירים עם עליית ה-SNR : בעוד אותה רשת MLP בת 169 פרמטרים משקמת ממסר אמין עד כדי מרחק של $1$–$1.5$ dB מגלאי Viterbi MLSE בעל ידיעת ערוץ מלאה, וללא כל ידע על משפחת הערוץ. התוצאה נתחמת מכל כיוון: היא לעולם אינה עולה על Viterbi מותאם; על שרשור מורכב של ISI, אי-ליניאריות ופאזה בלתי-ידועה היא נעה כ-$2$ dB מאחורי MLSE מבוסס-פיילוטים ללא כל מודל של הפגם; במשטר נטול-מידע-מוקדם היא משתווה למאזן עיוור (CMA) תוך הימנעות מאי-יציבות של MLSE מבוסס-החלטות; וסריקת תקציב פיילוטים מאתרת סף זיהוי חד הקובע מתי הממסר הלמוד, שאינו נדרש לפיילוטים כלל, עוקף עיבוד קלאסי מבוסס-פיילוטים. לבסוף, ניתוח מורכבות מראה שעלות הממסר הלמוד לכל סימבול קבועה בסדר האפנון ובזיכרון הערוץ (בעוד Viterbi גדל כ-$M^{L}$) ומהירה פי $30$–$90$ בזמן-ריצה. הממסר הלמוד אינו עולה על מקלט קלאסי מותאם, אך תופס נישה מוגדרת היטב של מיתון בלתי-תלוי-משפחה, נטול-זיהוי, וקבוע-מורכבות.

% אסטרטגיית הפריסה המומלצת היא ממסר Hybrid המשלב עיבוד למידה ב-SNR נמוך עם DF ב-SNR גבוה (169 פרמטרים, כ-0.7 KB, פחות מ-3 שניות אימון); עבור 16-QAM, ניסוח כמסווג דו-ממדי משותף; ובכל קישור הנושא פגם בלתי-ממודל (זיכרון, אי-ליניאריות), ממסר למוד מינימלי. התובנה המרכזית היא שממסרים למודים לעולם אינם עולים על האופטימום מודע-המודל, אך ייחודיים בעמידותם כאשר \emph{אין ידיעה איזה מודל חל} : ושמורכבות המודל חייבת להיות מותאמת למבנה המשימה, לא רק לגודלה.

תזה זו משווה בין אסטרטגיות ממסר קלאסיות ונלמדות לתקשורת דו־דילוגית. תצורה אחת קבועה לאורך המחקר המרכזי: קישור כניסה יחידה ויציאה יחידה (SISO) עם ממסר בודד, דעיכת ריילי מהירה בלתי־תלויה ושווה־התפלגות בשני הדילוגים, פס־בסיס מרוכב, אפנון QPSK מקודד־גריי ותמסורת בלתי־מקודדת, כאשר הממסר הוא המשתנה היחיד. שמונה אסטרטגיות מוערכות, מהגבר־והעבר (AF) ופענח־והעבר סימבולי (DF) ועד מודלים גנרטיביים ומודלי רצפים. התוצאות הן אומדני מונטה קרלו עם רווחי סמך של 95\%, מכוילים מול הביטוי הסגור להרכבה הדו־דילוגית.

בתרחיש זה, הממסרים הנלמדים הטובים ביותר עולים על AF ביחס אות־רעש (SNR) נמוך; ב־SNR בינוני וגבוה DF משתווה לכל ממסר או עולה עליו ללא עלות פרמטרים -- קביעה אמפירית על ממסר בלתי־מקודד וסימבולי־לפי־סימבול, ולא אופטימליות כללית. סריקת גדלים על פני כל הערוצים הנחקרים, בשלושה אתחולים לכל תצורה, מאתרת את הממסר הקטן ביותר שאינו כרוך בעלות מדידה, ומראה שזיכרון הערוץ, ולא מספר הפרמטרים, הוא הקובע אותו: בערוץ חסר־זיכרון ארבעה פרמטרים, יחידה נסתרת אחת ללא חלון, נמצאים במרחק 0.03 dB מן הבסיס הקלאסי, בעוד ערוצים בעלי שלושה מקדמים דורשים 73 עד 145 פרמטרים וחלון הפורש את ההפרעה. השגיאה אינה עולה שוב עם גידול הקיבולת, ולפיכך לא נצפתה התאמת־יתר בטווח הגדלים שנסרק. ממסר Hybrid העובר ל-DF מעל נקודת המעבר הנמדדת קרוב לאופטימלי ללא עלות נוספת, אם ה-SNR מוערך באמינות. מחקר מקודד מעלה את רף ההשוואה: קוד קונבולוציוני בקצב 1/2 עם ממסר DF מבוסס־בלוק, שונה מה־DF הסימבולי, עולה על DF בלתי־מקודד מ־8 dB ומעלה. שני הממסרים הנלמדים עוקפים אותו עד 20 dB (Mamba-S6 כבר ב־16): ממסר המפענח צורך את יתירות הקוד, ממסר המסיר רעש משאיר אותה שלמה. הסתגלות סכמת האפנון־והקידוד גורמת לסכמה להכריע את בחירת הממסר בסדר גודל -- אלא אם ההשהיה מחזירה את הפער.

התרומה העיקרית היא סולם של ארבע שכבות של הנחות ערוץ מוקלות בהדרגה. בשכבה 1 העיבוד הקלאסי מנצח: בערוץ חסר־זיכרון ותואם, החלטה קשה היא פונקציית הממסר האופטימלית לכל סימבול. שכבה 2 מוסיפה מסנן ISI תלת־מקדמי בלתי־ידוע, והממסרים הקלאסיים חסרי־הזיכרון נכשלים, כאשר שיעור השגיאה של DF עולה עם עוצמת השידור, בעוד ממסר נלמד ממוסגר־חלון בן 169 פרמטרים משקם את הקישור בעלות קבועה לכל סימבול, במרחק 1 עד 1.5~dB מגלאי Viterbi בעל עלות אקספוננציאלית. שכבה 3 מדרדרת את אומדן הערוץ: היתרון הקלאסי מחזיק עד כעשרים פיילוטים ואובד בעשרה. שכבה 4 מגרילה ערוץ חדש בכל בלוק, ושם הממסר הנלמד מוביל על פני המאזן העיוור הקלאסי הטוב ביותר.

התוצאות מתחמות משטרי הפעלה: עיבוד קלאסי כאשר מודל הערוץ המשוער תואם לערוץ, והממסר הנלמד כשיש לערוץ מבנה בלתי־ממודל, אומדנו מדורדר, או חסר לכל בלוק.



\end{flushright}
\end{justify}
\end{otherlanguage}
